%% ── 13. CONCLUSIONES ─────────────────────────────────────────────────────────
\chapter{Conclusiones}

El presente trabajo diseñó, implementó y validó ContractIA, un agente de
inteligencia artificial especializado en la auditoría automatizada de contratos
de concesión de infraestructura bajo el marco APP peruano. El sistema integra
técnicas de recuperación aumentada de generación (RAG híbrido), grafos de
conocimiento (GraphRAG), orquestación multi-agente y despliegue en nube
serverless, logrando un producto funcional accesible vía web y Telegram.

A continuación se presentan las conclusiones ordenadas por los tres ejes
del problema original: eficiencia operativa, calidad técnica del análisis
y viabilidad institucional.

\section*{Sobre la Eficiencia Operativa}

\textbf{Primera conclusión:}
El MVP demostró una reducción del tiempo de screening inicial de 320 horas
a $\approx 75$ minutos, equivalente a una reducción del \textbf{99.6\%} del
esfuerzo humano requerido para la revisión inicial de un contrato de
concesión de gran extensión. Este resultado supera el objetivo SMART
establecido ($\approx 1$ hora), valida la hipótesis central del proyecto
y confirma la viabilidad técnica del enfoque propuesto.

La reducción de tiempo no implica simplificación del análisis: el sistema
procesó \textbf{100\%} de las 41 secciones y 448 cláusulas del Contrato A, frente al $\approx 60\%$ de cobertura del proceso manual por muestreo. La cobertura total es una mejora estructural, no marginal, que permite detectar inconsistencias en secciones que históricamente no
eran revisadas en los plazos institucionales disponibles.

\textbf{Segunda conclusión:}
El costo operativo por análisis ($\approx \$50$--$\$100$ en tokens e infraestructura GCP) representa una reducción del \textbf{98.6\%} respecto al costo manual estimado ($\approx \$7{,}200$ en horas-hombre). A escala institucional, esto hace viable el análisis sistemático de todos los contratos de una cartera APP, no solo de los contratos seleccionados por muestreo.

\section*{Sobre la Calidad Técnica del Análisis}

\textbf{Tercera conclusión:}
La arquitectura multi-agente (Jurista/Auditor/Cronista) con RAG híbrido
(FAISS + BM25 + Cohere Reranker) y GraphRAG (NetworkX) detectó
\textbf{119 hallazgos} en el contrato piloto, con los siguientes indicadores
de calidad validados:

\begin{table}[H]
\centering
\caption{Resumen de indicadores de calidad del MVP}
\label{tab:resumen_calidad}
\begin{tabular}{p{4cm} p{3cm} p{3cm} p{3cm}}
\toprule
\textbf{Indicador} & \textbf{Objetivo} & \textbf{Resultado} & \textbf{Estado} \\
\midrule
Precisión semántica & $\geq 85\%$ & $\approx 87\%$ & \checkmark\ Alcanzado \\
Recall estimado     & $\geq 90\%$ & $\approx 94\%$ & \checkmark\ Alcanzado \\
Cobertura secciones & $100\%$     & $100\%$ (41/41) & \checkmark\ Alcanzado \\
Trazabilidad        & $100\%$     & $100\%$ (119/119) & \checkmark\ Alcanzado \\
Total hallazgos     & $> 30$      & 119 ($\times 3.9$) & \checkmark\ Superado \\
\bottomrule
\end{tabular}
\end{table}

La distribución uniforme de hallazgos entre los cuatro tipos
(\texttt{ERROR\_PLAZO} 29\%, \texttt{INCONSISTENCIA} 26\%,
\texttt{ERROR\_LOGICO} 23\%, \texttt{OMISION\_NORMATIVA} 23\%) confirma
que el sistema cubre de manera equilibrada las dimensiones temporal,
semántica, lógica y normativa del análisis contractual, sin sesgo hacia
un único tipo de inconsistencia.

\textbf{Cuarta conclusión:}
La evaluación comparativa entre Gemini~2.5~Pro y Gemini~3.1~Pro~Preview
reveló una \textbf{paradoja metodológica con valor científico}: el modelo
con grafo de conocimiento más rico (477 nodos, 617 relaciones vs.\ 208/258)
detectó significativamente \textit{menos} hallazgos (67 vs.\ 119), pero con
una distribución de severidad casi idéntica ($\approx 57\%$ ALTA en ambos).
Esta observación sugiere que un LLM más capaz usa el grafo como mecanismo
de filtrado de alta precisión, mientras que un modelo de menor capacidad
relacional lo usa como amplificador de cobertura.

Esta dicotomía cobertura--precisión entre modelos, mediada por la densidad
del grafo de conocimiento, constituye un \textbf{aporte metodológico} al
diseño de sistemas de auditoría multi-modelo: la elección del modelo no
puede hacerse solo por benchmarks genéricos, sino que debe considerar el
trade-off cobertura/precisión en el dominio específico.

\textbf{Quinta conclusión:}
El \textbf{100\% de trazabilidad} logrado --- cada hallazgo con ubicación
exacta, evidencia textual, referencia normativa y razonamiento explícito ---
constituye el diferenciador institucional más importante del sistema.
Sin trazabilidad completa, los hallazgos generados por IA no son accionables
en procesos formales de auditoría contractual. Esta característica aborda
directamente uno de los principales puntos de dolor reportados por
especialistas de ProInversión: la ausencia de documentación estructurada
en la revisión manual.

\section*{Sobre la Viabilidad Institucional}

\textbf{Sexta conclusión:}
El despliegue exitoso del MVP en producción sobre Google Cloud Run, con
frontend Next.js accesible en \texttt{contractia.pe}, bot Telegram operativo,
sistema de cola por usuario, notificación por correo y CI/CD automatizado,
demostró que la solución no es solo un prototipo académico sino un sistema
\textbf{operativo, mantenible y escalable}. La arquitectura serverless
elimina los costos fijos de infraestructura y permite escalar a cero cuando
no hay análisis en curso, haciendo viable la operación institucional con
presupuestos de TI moderados.

\textbf{Séptima conclusión:}
El enfoque iterativo de desarrollo (15 versiones de notebook + 4 versiones
de plataforma) y el uso del mismo contrato piloto como benchmark constante
validaron que el ciclo \textit{build--measure--learn} es aplicable y efectivo
en proyectos de IA aplicada. La metodología permitió reducir el riesgo técnico
acumulativo, detectar problemas de producción tempranamente y comunicar
el progreso con evidencia cuantitativa en cada iteración.

\textbf{Octava conclusión:}
ContractIA no pretende reemplazar al especialista legal ni al auditor
contractual. Su rol definido es el de \textbf{herramienta de screening
previo}: acelerar la fase de identificación inicial de inconsistencias,
ampliar la cobertura al 100\% del documento y generar un insumo estructurado
y trazable que el especialista puede revisar, validar y convertir en
observaciones formales. El $\approx 13\%$ de falsos positivos estimado
confirma que el juicio experto sigue siendo indispensable para la toma de
decisiones finales.

\section*{Contribución al Campo}

ContractIA constituye, en el contexto peruano, uno de los primeros sistemas
operativos de auditoría contractual basado en LLMs con los siguientes aportes
diferenciados:

\begin{enumerate}
    \item \textbf{Arquitectura multi-agente especializada por tipo de
    inconsistencia} en el dominio de contratos de concesión APP, con esquemas
    de salida estructurados y trazabilidad completa.
    \item \textbf{Pipeline RAG híbrido} (vectorial + léxico + cross-encoder)
    adaptado al vocabulario técnico-legal del marco normativo peruano de APPs.
    \item \textbf{Integración de GraphRAG} para detección de inconsistencias
    entre cláusulas no adyacentes, habilitando un tipo de hallazgo
    estructuralmente imposible en revisiones lineales.
    \item \textbf{Demostración empírica} de la paradoja cobertura--precisión
    en sistemas GraphRAG con modelos LLM de distinta capacidad relacional.
    \item \textbf{Plataforma institucional desplegada en producción},
    incluyendo autenticación, roles, cola de trabajos, notificaciones y
    CI/CD, lista para pilotaje con usuarios reales de ProInversión.
\end{enumerate}

\section*{Perspectiva Final}

Los resultados obtenidos permiten afirmar que la integración de LLMs,
RAG híbrido y GraphRAG es una estrategia efectiva y viable para automatizar
el análisis de documentos contractuales complejos. La reducción del 99.6\%
en tiempo, la cobertura total y la multiplicación por 4.8 en hallazgos
detectados demuestran un impacto operativo concreto y medible.

ContractIA se posiciona como un \textbf{habilitador clave para la
transformación digital del proceso de auditoría contractual en el Perú},
con potencial de extensión a otros tipos de contratos, entidades del
sector público y, en el largo plazo, a otras jurisdicciones
latinoamericanas con marcos APP similares.
